% =====================================================================
\section{Experimental Environment}

\subsection{A Literature-Calibrated Synthetic Environment}
The environment is not a replica of the operating data of a particular
terminal but a \emph{literature-calibrated synthetic environment} that
combines the scales and operating ranges used in existing container
terminal studies. Table~\ref{tab:env} separates values taken directly
from prior work from values this study fixed within those ranges.

\begin{table}[t]
\centering
\caption{Environment components: basis in prior work and the value used here}
\label{tab:env}
\small
\begin{tabular}{@{}lll@{}}
\toprule
Component & Basis in prior work & Value used here \\
\midrule
Multiple blocks & 16 and 20 blocks~\cite{ref10,ref17} & 21 blocks \\
Multiple yard cranes & Several cranes per block~\cite{ref5,ref18} & 2 per block \\
Long-run discrete event & Multi-week continuous runs~\cite{ref13} & 30 days, state carried \\
Time-of-day arrivals & Hourly probability, turn time~\cite{ref19,ref20} & Base plus three peaks \\
Vessel class, call size & Three classes, differing calls~\cite{ref17,ref21} & Small/medium/large \\
Crane service time & Triangular handling range~\cite{ref13} & 180\,s within the range \\
\bottomrule
\end{tabular}
\end{table}

The 21 blocks, the exact demand levels, and the arrival-curve
coefficients are therefore design values of this study. The citations
establish that the chosen scale and variation lie within the modelling
range of existing work; they do not indicate operating data from a
specific port.

\subsection{Terminal and Time Settings}
The yard has 21 blocks with two yard cranes each, 42 in total. Block
capacity is 1{,}440 slots with an initial occupancy of 65\%. The layout
is one-dimensional and horizontal, with the gate and the quay on
opposite sides. The policy decision interval is 60\,s. Operation runs
for 30 continuous days, and the middle 28 days are compared after
excluding the first and last day. The crane work sequence gives priority
to actual service work and, within it, to the job with the shortest
expected processing time.

A day is not a point at which the terminal stops but a cut at which cost
and intermediate results are separated. Yard inventory, queues, crane
state, and unfinished jobs carry over to the next day. The first day
connects the initial state to the continuous flow, and the last day
preserves the residual effect of preceding work.

\subsection{External Truck Demand and the Arrival Process}
Daily demand is drawn from 3{,}500 (30\%), 5{,}000 (30\%), 7{,}500
(25\%), 12{,}500 (10\%), and 15{,}000 (5\%) trucks. The first three
represent smooth, normal, and high demand; the last two represent
congested and heavily congested conditions. The time-of-day arrival
density $f(t)$ is a mixture of a uniform base over 24 hours and normal
peaks:

\begin{equation}
f(t)=0.38\,U(0,24)+0.62\sum_{m=1}^{3}w_m
\mathcal N_{[0,24)}(t;\mu_m,\sigma_m^2),
\end{equation}
\begin{equation}
(\mu_m,\sigma_m,w_m)\in
\{(10,1.5,0.317),(15,2.5,0.633),(21,1.0,0.050)\}.
\end{equation}

Time is measured in hours and the truncated density is renormalised so
that it integrates to one over a day. The time-of-day arrival model and
its effect on turn time follow~\cite{ref19,ref20}; the mixture weights
and peak centres are design values chosen to compare congestion
sensitivity. Fig.~\ref{fig:arrival} shows the resulting curve.

\begin{figure}[t]
\centering
\includegraphics[width=0.86\textwidth]{fig2-arrival.pdf}
\caption{Time-of-day truck arrival density: a uniform base of 38\% over
24 hours plus three normal peaks.}
\label{fig:arrival}
\end{figure}

\subsection{Vessels and Quay Operations}
Vessels are divided into three classes and the call volume is generated
as a fixed fraction of vessel capacity (Table~\ref{tab:vessel}). Call
volume is $\text{TEU}\times U(0.15,0.30)$. The default daily composition
is two small vessels and one medium vessel, with one large vessel added
with probability 0.30. The throughput of one quay crane is set to
27.5\,move/h, which lies within the 24, 30, and 36\,move/h range used in
the literature~\cite{ref17}. Blocking across quay cranes is summed and
divided by the number assigned, giving idle time per vessel.

\begin{table}[t]
\centering
\caption{Vessel classes and call volumes}
\label{tab:vessel}
\small
\begin{tabular}{@{}lrrrrr@{}}
\toprule
Class & GT & TEU & Call volume & Quay cranes & Idle cost/hour \\
\midrule
Small & 50,000 & 3,000 & 450--900 move & 2 & KRW 149,500 \\
Medium & 100,000 & 7,500 & 1,125--2,250 move & 4 & KRW 299,000 \\
Large & 150,000 & 14,000 & 2,100--4,200 move & 6 & KRW 448,500 \\
\bottomrule
\end{tabular}
\end{table}

\subsection{Units of Yard Crane Capacity}
The reference service time of one crane is 180\,s per move, so the
nominal total move capacity is

\begin{equation}
\mu_{YC}=21\times2\times\frac{3600}{180}=840\ \text{move/h}.
\end{equation}

A move is a container movement and is not a truck. One truck job may
require several moves and rehandles, and the realised service time
varies with position and job type. We follow the guidance that equipment
capacity must be reported together with move type, automation delay, and
handling-vehicle conditions~\cite{ref6}.

\subsection{Cost Function}
The total cost over the interval $[d,d+1)$ is the sum of four terms:

\begin{equation}
\Phi_d=C_{wait,d}+C_{move,d}+C_{rehandle,d}+C_{vessel,d}.
\end{equation}

Truck dwell costs KRW 40{,}000 per hour, with the same rate applied
again to any excess beyond one hour; the safe trucking freight rate is
used as the monetary anchor and the hourly conversion is a study
value~\cite{ref12}. Additional crane movement costs KRW 60{,}000 per
hour of extra work and a rehandle costs KRW 2{,}000, both design values.
Vessel idling costs KRW 2.99 per GT-hour, converted from the excess
berthing rate of KRW 29.9 per 10\,GT-hour~\cite{ref22}.

With $h_i$ the terminal dwell time of truck $i$, $m$ the additional
crane working time, $R$ the number of rehandles, and $T_v$ the
policy-related idle time of vessel $v$:

\begin{equation}
C_{wait}=\sum_i 40{,}000\{h_i+\max(0,h_i-1)\},\quad
C_{move}=60{,}000\,m,\quad C_{rehandle}=2{,}000\,R,
\end{equation}
\begin{equation}
C_{vessel}=\sum_v 2.99\,GT_v\,T_v.
\end{equation}

Structural vessel time that the reallocation policy cannot change, such
as berthing and quarantine, is separated as a diagnostic value. All four
terms are computed as differences of cumulative values at the same day
boundary.
